\chapter{ตารางค่าดัชนีหักเหและความหนาแน่นของสารละลายมาตรฐาน}
\label{app:icumsa}

ตารางที่ \ref{tab:appendix_icumsa} แสดงค่าดัชนีหักเหของแสง (Refractive Index, $n$) และค่าความหนาแน่นของสารละลาย ($\rho$) ที่ระดับความเข้มข้นสารละลายของแข็งและน้ำตาลต่าง ๆ ($^\circ\text{Bx}$) ที่อุณหภูมิ 20 องศาเซลเซียส ตามมาตรฐานสากล ICUMSA และแบบจำลอง Lorentz-Lorenz

\begin{table}[H]
\centering
\caption{ค่ามาตรฐานดัชนีหักเหและความหนาแน่นของสารละลายตามเกณฑ์ ICUMSA}
\label{tab:appendix_icumsa}
\begin{tabularx}{\textwidth}{YYYY}
\toprule
\textbf{ความเข้มข้น ($^\circ\text{Bx}$)} & \textbf{ดัชนีหักเห ($n$)} & \textbf{ความหนาแน่น ($\text{g/cm}^3$)} & \textbf{การนำไปใช้เทียบเคียง} \\
\midrule
0.0 (น้ำกลั่นบริสุทธิ์) & 1.3330 & 0.9982 & จุดอ้างอิงศูนย์ (Blank Reference) \\
5.0 & 1.3403 & 1.0177 & สารละลายธาตุอาหารดินเจือจาง \\
10.0 & 1.3478 & 1.0381 & สารละลายน้ำสกัดผลไม้หรือน้ำหมัก \\
15.0 & 1.3557 & 1.0592 & สารละลายน้ำตาลกลูโคสเข้มข้น \\
20.0 & 1.3639 & 1.0810 & สารละลายปุ๋ยน้ำเข้มข้น \\
\bottomrule
\end{tabularx}
\end{table}

\vspace{0.8cm}
\clearpage
\chapter{ข้อมูลจำเพาะเชิงเทคนิคและสถาปัตยกรรมฮาร์ดแวร์ Wio Terminal}
\label{app:specs}

ตารางที่ \ref{tab:appendix_specs} สรุปข้อมูลจำเพาะเชิงเทคนิคและสถาปัตยกรรมระบบประมวลผลของเครื่องสเปกโทรโฟโตมิเตอร์ Wio Terminal ATSAMD51

\begin{table}[H]
\centering
\caption{ข้อมูลจำเพาะเชิงเทคนิคของระบบประมวลผล Wio Terminal ATSAMD51}
\label{tab:appendix_specs}
\begin{tabularx}{\textwidth}{p{4.5cm}X}
\toprule
\textbf{พารามิเตอร์ทางเทคนิค} & \textbf{รายละเอียดคุณลักษณะ} \\
\midrule
ไมโครคอนโทรลเลอร์ & Microchip ATSAMD51P19A (ARM Cortex-M4F 120 MHz) \\
หน่วยประมวลผลทศนิยม & Hardware Floating Point Unit (FPU Single-Precision 32-bit) \\
หน่วยความจำโปรแกรมแฟลช & 512 KB Flash Memory (ใช้งานจริงประมาณ 103 KB คิดเป็น 20\%) \\
หน่วยความจำแรมภายใน & 192 KB SRAM \\
จอแสดงผลข้อมูล & TFT LCD ขนาด 2.4 นิ้ว ความละเอียด 320x240 พิกเซล \\
โมดูลเซนเซอร์ตรวจวัดสี & Adafruit TCS34725 (RGBC Photodiode Array พร้อม IR Filter) \\
โมดูลกำเนิดแสงสเปกตรัม & Grove WS2812B NeoPixel (ควบคุมความยาวคลื่นแม่นยำ 5 แถบแสง) \\
ระบบจัดเก็บข้อมูลภายนอก & microSD Card Slot รองรับระบบไฟล์ FAT32 (Quadruple Logging) \\
\bottomrule
\end{tabularx}
\end{table}
